% problem_id: S001-lemma-canonical-class-obstruction
% source_id: S001 (Stacks Project)
% source_file: defos.tex
% statement_locator: defos.tex lines 5357--5372
% proof_locator: defos.tex lines 5374--5495
% env: lemma
% id_note: label 在本来源内唯一
% pairing: tight (verified)
% status: complete (statement + author proof verbatim + reason + locators)
%
\begin{lemma}
\label{lemma-canonical-class-obstruction}
Let $\mathcal{C}$ be a site. Let $\mathcal{O} \to \mathcal{O}_0$
be a surjection of sheaves of rings whose kernel is an ideal sheaf
$\mathcal{I}$ of square zero. For every object
$K_0$ in $D^-(\mathcal{O}_0)$ the following are equivalent
\begin{enumerate}
\item  the class
$\omega(K_0) \in
\Ext^2_{\mathcal{O}_0}(K_0, K_0 \otimes_{\mathcal{O}_0} \mathcal{I})$
constructed in Lemma \ref{lemma-canonical-class} is zero,
\item there exists $K \in D^-(\mathcal{O})$ with
$K \otimes_\mathcal{O}^\mathbf{L} \mathcal{O}_0 = K_0$
in $D(\mathcal{O}_0)$.
\end{enumerate}
\end{lemma}

\begin{proof}
Let $K$ be as in (2). Then we can represent $K$ by a bounded above
complex $\mathcal{F}^\bullet$ of flat $\mathcal{O}$-modules.
Then $\mathcal{F}_0^\bullet =
\mathcal{F}^\bullet \otimes_{\mathcal{O}} \mathcal{O}_0$
represents $K_0$ in $D(\mathcal{O}_0)$.
Since $d_{\mathcal{F}^\bullet} \circ d_{\mathcal{F}^\bullet} = 0$
as $\mathcal{F}^\bullet$ is a complex, we see from the very construction
of $\omega(K_0)$ that it is zero.

\medskip\noindent
Assume (1). Let $\mathcal{F}^n$, $d : \mathcal{F}^n \to \mathcal{F}^{n + 1}$
be as in the construction of $\omega(K_0)$. The nullity of
$\omega(K_0)$ implies that the map
$$
\omega = d \circ d : \mathcal{F}_0^\bullet
\longrightarrow
(\mathcal{F}_0^\bullet \otimes_{\mathcal{O}_0} \mathcal{I})[2]
$$
is zero in $D(\mathcal{O}_0)$. By definition of the derived category
as the localization of the homotopy category of complexes
of $\mathcal{O}_0$-modules, there exists a quasi-isomorphism
$\alpha : \mathcal{G}_0^\bullet \to \mathcal{F}_0^\bullet$
such that there exist $\mathcal{O}_0$-modules maps
$h^n : \mathcal{G}_0^n \to
\mathcal{F}_0^{n + 1} \otimes_\mathcal{O} \mathcal{I}$
with
$$
\omega \circ \alpha =
d_{\mathcal{F}_0^\bullet \otimes \mathcal{I}} \circ h +
h \circ d_{\mathcal{G}_0^\bullet}
$$
We set
$$
\mathcal{H}^n = \mathcal{F}^n \times_{\mathcal{F}^n_0} \mathcal{G}_0^n
$$
and we define
$$
d' : \mathcal{H}^n \longrightarrow \mathcal{H}^{n + 1},\quad
(f^n, g_0^n) \longmapsto (d(f^n) - h^n(g_0^n), d(g_0^n))
$$
with obvious notation using that
$\mathcal{F}_0^{n + 1} \otimes_{\mathcal{O}_0} \mathcal{I} =
\mathcal{F}^{n + 1} \otimes_\mathcal{O} \mathcal{I} =
\mathcal{I}\mathcal{F}^{n + 1} \subset \mathcal{F}^{n + 1}$.
Then one checks $d' \circ d' = 0$ by our choice of $h^n$
and definition of $\omega$.
Hence $\mathcal{H}^\bullet$ defines an object in $D(\mathcal{O})$.
On the other hand, there is a short exact sequence of complexes
of $\mathcal{O}$-modules
$$
0 \to \mathcal{F}_0^\bullet \otimes_{\mathcal{O}_0} \mathcal{I} \to
\mathcal{H}^\bullet \to \mathcal{G}_0^\bullet \to 0
$$
We still have to show that
$\mathcal{H}^\bullet \otimes_\mathcal{O}^\mathbf{L} \mathcal{O}_0$
is isomorphic to $K_0$.
Choose a quasi-isomorphism
$\mathcal{E}^\bullet \to \mathcal{H}^\bullet$
where $\mathcal{E}^\bullet$ is a bounded above complex of flat
$\mathcal{O}$-modules. We obtain a commutative diagram
$$
\xymatrix{
0 \ar[r] &
\mathcal{E}^\bullet \otimes_\mathcal{O} \mathcal{I} \ar[d]^\beta \ar[r] &
\mathcal{E}^\bullet \ar[d]^\gamma \ar[r] &
\mathcal{E}_0^\bullet \ar[d]^\delta \ar[r] &
0 \\
0 \ar[r] &
\mathcal{F}_0^\bullet \otimes_{\mathcal{O}_0} \mathcal{I} \ar[r] &
\mathcal{H}^\bullet \ar[r] &
\mathcal{G}_0^\bullet \ar[r] &
0
}
$$
We claim that $\delta$ is a quasi-isomorphism. Since $H^i(\delta)$
is an isomorphism for $i \gg 0$, we can use descending induction
on $n$ such that $H^i(\delta)$ is an isomorphism for $i \geq n$.
Observe that
$\mathcal{E}^\bullet \otimes_\mathcal{O} \mathcal{I}$
represents
$\mathcal{E}_0^\bullet \otimes_{\mathcal{O}_0}^\mathbf{L} \mathcal{I}$,
that
$\mathcal{F}_0^\bullet \otimes_{\mathcal{O}_0} \mathcal{I}$
represents
$\mathcal{G}_0^\bullet \otimes_{\mathcal{O}_0}^\mathbf{L} \mathcal{I}$,
and that
$\beta = \delta \otimes_{\mathcal{O}_0}^\mathbf{L} \text{id}_\mathcal{I}$
as maps in $D(\mathcal{O}_0)$. This is true because
$\beta =
(\alpha \otimes \text{id}_\mathcal{I})
\circ
(\delta \otimes \text{id}_\mathcal{I})$.
Suppose that $H^i(\delta)$ is an isomorphism in degrees $\geq n$.
Then the same is true for $\beta$ by what we just said
and Lemma \ref{lemma-induced-map}.
Then we can look at the diagram
$$
\xymatrix{
H^{n - 1}(\mathcal{E}^\bullet \otimes_\mathcal{O} \mathcal{I})
\ar[r] \ar[d]^{H^{n - 1}(\beta)} &
H^{n - 1}(\mathcal{E}^\bullet) \ar[r] \ar[d] &
H^{n - 1}(\mathcal{E}_0^\bullet) \ar[r] \ar[d]^{H^{n - 1}(\delta)} &
H^n(\mathcal{E}^\bullet \otimes_\mathcal{O} \mathcal{I})
\ar[r] \ar[d]^{H^n(\beta)} &
H^n(\mathcal{E}^\bullet) \ar[d] \\
H^{n - 1}(\mathcal{F}_0^\bullet \otimes_\mathcal{O} \mathcal{I}) \ar[r] &
H^{n - 1}(\mathcal{H}^\bullet) \ar[r] &
H^{n - 1}(\mathcal{G}_0^\bullet) \ar[r] &
H^n(\mathcal{F}_0^\bullet \otimes_\mathcal{O} \mathcal{I}) \ar[r] &
H^n(\mathcal{H}^\bullet)
}
$$
Using Homology, Lemma \ref{homology-lemma-four-lemma}
we see that $H^{n - 1}(\delta)$ is surjective.
This in turn implies that $H^{n - 1}(\beta)$ is surjective
by Lemma \ref{lemma-induced-map}.
Using Homology, Lemma \ref{homology-lemma-four-lemma}
again we see that $H^{n - 1}(\delta)$ is an isomorphism.
The claim holds by induction, so $\delta$ is a quasi-isomorphism
which is what we wanted to show.
\end{proof}
